% Vietnamese translation for OI Public Library
% Source-File: IOI中国国家候选队论文/2003/伍昱/伍昱.ppt
\documentclass[11pt]{article}
\usepackage{oipl-vn}

\title{Bản trình chiếu: Giải bài toán 2-SAT bằng tính đối xứng}
\author{Wu Yu}
\date{2003}

\begin{document}
\maketitle

\origtitle{Solving 2-SAT with symmetry}
\sourcepath{IOI China candidate team papers / 2003 / Wu Yu / Wu Yu.ppt}

\section{Bài toán xuất phát}

Bản trình chiếu mở đầu bằng bài \textit{Peaceful Commission}. Có \(n\)
đảng, mỗi đảng có đúng hai đại biểu, đánh số \(2a-1\) và \(2a\) cho đảng
thứ \(a\). Ta phải chọn đúng một đại biểu của mỗi đảng để lập ủy ban hòa
bình. Một số cặp đại biểu không thân thiện với nhau; hai người trong cùng
một cặp như vậy không được cùng xuất hiện trong ủy ban. Với
\(1\le n\le 8000\) và \(m\le 20000\), cách thử từng tổ hợp là hoàn toàn
không khả thi.

Nội dung chính của slide là biến bài toán này thành một bài 2-SAT có cấu
trúc rất đều. Mỗi lựa chọn \(A_i\) luôn có lựa chọn đối ứng \(A_i'\) trong
cùng một nhóm, và mỗi nghiệm phải chọn đúng một trong hai lựa chọn đó.

\section{Cung kéo theo và tính đối xứng}

Nếu hai lựa chọn \(A_i\) và \(A_j\) không tương thích, chọn một lựa chọn sẽ
buộc phải chọn lựa chọn đối ứng của lựa chọn kia:
\[
  A_i \to A_j', \qquad A_j \to A_i'.
\]
Đây là đồ thị kéo theo của công thức 2-SAT. Điểm quan trọng được các slide
nhấn mạnh là mọi cung đều có cung đối xứng: nếu có \(A_i\to A_j\), thì có
\(A_j'\to A_i'\). Tính chất này vẫn đúng cho cả đường đi, không chỉ cho
một cung đơn lẻ; từ \(A_i\to^* A_k\) suy ra \(A_k'\to^* A_i'\).

Vì vậy, khi ta chọn một đỉnh trong đồ thị kéo theo, tất cả các đỉnh đi được
từ nó cũng bị buộc chọn. Đồng thời các đỉnh đối ứng với chúng bị loại. Nếu
một lựa chọn và lựa chọn đối ứng của nó cùng bị buộc, bài toán vô nghiệm.

\section{Thành phần liên thông mạnh}

Bản trình chiếu tiếp tục gom các đỉnh nằm trong cùng một chu trình có hướng.
Nếu một đỉnh trong chu trình được chọn, mọi đỉnh còn lại trong chu trình
cũng bị buộc chọn, nên chúng có thể được co thành một thành phần liên thông
mạnh. Sau khi co tất cả các thành phần, ta thu được một DAG.

Điều kiện vô nghiệm trở nên rất rõ:
\[
  A_i \text{ và } A_i' \text{ thuộc cùng một thành phần liên thông mạnh}
  \quad\Longrightarrow\quad \text{vô nghiệm}.
\]
Ngược lại, nếu không có cặp đối ứng nào rơi vào cùng một thành phần, tính
đối xứng của đồ thị co bảo đảm có thể chọn nhất quán một phía trong từng
cặp thành phần đối ứng.

\section{Cách xây dựng nghiệm}

Một cách trực tiếp là chọn một thành phần chưa quyết định, lấy nó cùng mọi
hậu duệ trong DAG, rồi loại thành phần đối ứng cùng mọi tiền bối của thành
phần đối ứng. Nhờ tính đối xứng, hai tập này khớp nhau đúng theo quan hệ
đối ứng, nên thao tác không tạo mâu thuẫn mới nếu bài toán đã qua kiểm tra
thành phần liên thông mạnh.

Để thực hiện gọn hơn, các slide dùng thứ tự topo từ đáy lên. Khi duyệt DAG
theo thứ tự này, nếu gặp một thành phần chưa được quyết định thì chọn nó và
đánh dấu thành phần đối ứng là không chọn. Các hậu duệ đã được xử lý trước
đó, nên không cần lan truyền lại toàn bộ mỗi lần.

\section{Thuật toán tổng quát}

Quy trình hoàn chỉnh gồm các bước:
\begin{enumerate}
  \item Dựng đồ thị kéo theo từ các cặp không tương thích.
  \item Tìm các thành phần liên thông mạnh.
  \item Nếu một lựa chọn và lựa chọn đối ứng của nó cùng thuộc một thành
  phần, kết luận vô nghiệm.
  \item Co đồ thị thành DAG và lấy thứ tự topo từ đáy lên.
  \item Duyệt các thành phần theo thứ tự đó; thành phần nào chưa quyết định
  thì chọn nó và loại thành phần đối ứng.
  \item Từ các thành phần được chọn, xuất một đại biểu cho mỗi đảng.
\end{enumerate}

Nếu dùng Tarjan hoặc Kosaraju để tìm thành phần liên thông mạnh, mọi bước
sau khi dựng đồ thị đều tuyến tính theo số đỉnh và số cung. Với bài
\textit{Peaceful Commission}, số cung tỉ lệ với số cặp không thân thiện, nên
độ phức tạp là \(O(n+m)\).

\section{Thông điệp của bài trình bày}

Mục đích của bộ slide không chỉ là nhắc lại thuật toán 2-SAT chuẩn, mà còn
giải thích vì sao tính đối xứng của đồ thị kéo theo làm thuật toán đúng. Khi
nhận ra mỗi cung, mỗi đường đi và mỗi thành phần đều có đối ứng, ta không
cần thử từng lựa chọn một cách mù quáng nữa; việc co SCC và duyệt topo cho
ta một cách chọn nghiệm hiệu quả và có chứng minh rõ ràng.

\end{document}
